% preamble.tex - LaTeX-Einstellungen für PDF-Export
% TH Köln Design - Roman Bartnik 2026
% Eingebunden über _quarto.yml (format: pdf: include-in-header)

% Satzspiegel-Feinheiten für 17 x 24 cm (KOMA-Script)
% (Satzspiegel kommt aus geometry in _quarto.yml, daher kein DIV=calc)
\usepackage{microtype}
\usepackage{xurl}  % URLs dürfen an jeder Stelle umbrechen (lange Links in schmalen Tabellenspalten)
\raggedbottom

% Tabellen: kleinere Schrift (9 pt) und Flattersatz mit Silbentrennung in den Zellen.
% Pandoc setzt Zellen mit \raggedright, was die Trennung praktisch abschaltet; \RaggedRight (ragged2e) trennt weiter.
\usepackage{etoolbox}
\usepackage{needspace}  % Titelzeile eines Prompts nicht allein am Seitenende (_filters/druck-satz.lua)
% Zeichen, die TeX Gyre Pagella nicht enthält (sonst leere Kästchen im Druck)
\usepackage{newunicodechar}
\usepackage{pifont}
\newunicodechar{★}{\ding{72}}
\newunicodechar{☆}{\ding{73}}
\newunicodechar{₂}{\textsubscript{2}}
% Bildunterschriften aus Teilbildern: Rest nach dem Kurztitel, entfällt im
% Abbildungsverzeichnis (_filters/druck-satz.lua)
\DeclareRobustCommand{\druckrest}[1]{#1}
\AtBeginDocument{\addtocontents{lof}{\protect\renewcommand{\protect\druckrest}[1]{}}}
% Kurztitel mit Fußnote (seit 14.11.2026): In Kästen setzt Quarto die Fußnote einer
% Abbildungs- oder Tabellenunterschrift als \footnotemark ohne Kurztitel fürs Verzeichnis.
% Robust gemacht, gelangt das Zeichen unbeschadet in .lof/.lot; dort schaltet es der
% folgende Eintrag stumm.
\MakeRobust{\footnotemark}
\AtBeginDocument{%
  \addtocontents{lof}{\protect\renewcommand{\protect\footnotemark}[1][]{}}%
  \addtocontents{lot}{\protect\renewcommand{\protect\footnotemark}[1][]{}}}
\usepackage{ragged2e}
\newcommand{\tabellenschrift}{\small}  % je Tabelle übersteuerbar, z. B. \renewcommand{\tabellenschrift}{\footnotesize} in einer Gruppe
\AtBeginEnvironment{longtable}{\tabellenschrift\let\raggedright\RaggedRight}  % kein vertikales Strecken, wenn ein Kasten nicht mehr auf die Seite passt

% Tabellen in Kästen (Quarto-Callouts = tcolorbox) und Tabellen ohne Überschrift.
% (1) longtable gibt Fuß und Schlussfuß nur in seiner Ausgaberoutine aus, und die läuft
%     in einem Kasten nie. Pandoc setzt dort wegen cap-location: bottom die Schlusslinie,
%     die Überschrift und das Label hinein; im Kasten fehlten alle drei (Befund 01.11.2026:
%     Tabelle 2.3 und 3.2). Abhilfe: Steht die Tabelle in einer Box (\ifinner), hängt
%     \endlongtable den Schlussfuß selbst an.
% (2) longtable zählt jede Tabelle, auch ohne Überschrift; sieben unbeschriftete Tabellen
%     in Kästen des damaligen Anhangs D erzeugten die Lücke D.2 -> D.10. Abhilfe: Tabellen ohne
%     Überschrift geben ihre Nummer zurück. Damit die PDF-Sprungmarken eindeutig bleiben,
%     enthält \theHtable die laufende longtable-Nummer.
\makeatletter
\newif\ifLT@mitueberschrift
\AtBeginEnvironment{longtable}{\global\LT@mitueberschriftfalse}
\AtBeginDocument{%
  \patchcmd{\LT@caption}{\noalign\bgroup}{\noalign\bgroup\global\LT@mitueberschrifttrue}{}%
    {\PackageWarningNoLine{preamble}{Patch LT@caption fehlgeschlagen}}%
  \patchcmd{\endlongtable}
    {\endgraf\penalty -\LT@end@pen}
    {\ifinner\nobreak\box\ifvoid\LT@lastfoot\LT@foot\else\LT@lastfoot\fi\fi
     \ifLT@mitueberschrift\else\addtocounter{table}{-1}\fi
     \endgraf\penalty -\LT@end@pen}
    {}{\PackageWarningNoLine{preamble}{Patch endlongtable fehlgeschlagen}}%
  \@ifundefined{theHtable}{}{\renewcommand*{\theHtable}{\arabic{LT@tables}.\theHchapter.\arabic{table}}}%
}
\makeatother

\usepackage{booktabs}
\usepackage{xcolor}
\usepackage{tcolorbox}
\tcbuselibrary{skins,breakable}  % Quarto-Callouts brauchen diese Bibliotheken
% fontenc entfällt: XeLaTeX nutzt fontspec (siehe _quarto.yml, mainfont)

% ===== TH KÖLN FARBPALETTE =====
\definecolor{thred}{HTML}{C81E0F}
\definecolor{thmagenta}{HTML}{B43092}
\definecolor{thorange}{HTML}{EA5A00}
\definecolor{thgray}{HTML}{6B6B73}
\definecolor{thlightgray}{HTML}{F5F5F7}

% Alte Farben (für Kompatibilität mit bestehenden Callouts)
\definecolor{dangcol}{RGB}{152,62,130}
\definecolor{warncol}{RGB}{245,220,112}
\definecolor{infocol}{RGB}{70,122,172}
\definecolor{trycol}{RGB}{97,88,156}

% ===== CAPTION-STYLING =====
\usepackage{caption}
\captionsetup{
  labelfont={bf},
  textfont={},
  font=small,
  labelsep=period,
  justification=raggedright,
  singlelinecheck=false
}
% Grauer Caption-Text
\captionsetup{textfont={color=thgray}}

% ===== CODEBLÖCKE (PROMPTS) =====
% Lange Prompt-Zeilen umbrechen statt rechts abzuschneiden. Quarto definiert die
% Umgebung "Highlighting" ohne breaklines; diese Datei wird danach eingebunden
% und überschreibt die Definition. Plain-Codeblöcke (ohne Sprache) laufen über
% "verbatim" und werden ebenfalls umbrechend gesetzt.
\usepackage{fvextra}
\DefineVerbatimEnvironment{Highlighting}{Verbatim}{commandchars=\\\{\},breaklines,breakanywhere=false,breaksymbolleft={},breakindent=1em,fontsize=\small}
\RecustomVerbatimEnvironment{verbatim}{Verbatim}{breaklines,breakanywhere=false,breaksymbolleft={},breakindent=1em,fontsize=\small}

% ===== TABELLEN =====
\usepackage{colortbl}
% Tabellenköpfe mit Hintergrund
\newcommand{\thheadrow}{\rowcolor{thlightgray}}

% ===== HYPERLINKS =====
% Linkfarben werden in _quarto.yml gesetzt (linkcolor/urlcolor/citecolor = TH-Farben)

% ===== CALLOUT-BOXEN =====
\newtcolorbox{dangerous}{
  colback=dangcol!10,
  colframe=dangcol,
  coltext=black,
  boxsep=5pt,
  arc=4pt
}

\newtcolorbox{warning}{
  colback=warncol!10,
  colframe=warncol,
  coltext=black,
  boxsep=5pt,
  arc=4pt
}

\newtcolorbox{info}{
  colback=infocol!10,
  colframe=infocol,
  coltext=black,
  boxsep=5pt,
  arc=4pt
}

\newtcolorbox{try}{
  colback=trycol!10,
  colframe=trycol,
  coltext=black,
  boxsep=5pt,
  arc=4pt
}

% Neue TH-Köln-Boxen
\newtcolorbox{thbox}{
  colback=white,
  colframe=thmagenta,
  coltext=black,
  boxsep=5pt,
  arc=0pt,
  left=3pt,
  leftrule=4pt,
  rightrule=0.5pt,
  toprule=0.5pt,
  bottomrule=0.5pt
}

\newtcolorbox{thwarning}{
  colback=white,
  colframe=thorange,
  coltext=black,
  boxsep=5pt,
  arc=0pt,
  left=3pt,
  leftrule=4pt,
  rightrule=0.5pt,
  toprule=0.5pt,
  bottomrule=0.5pt
}

% ------------------------------------------------------------------
% Impressum auf der Rückseite der Titelseite (KOMA-Script \lowertitleback).
% Angaben in [eckigen Klammern] sind Platzhalter, die vor der Druckfreigabe
% zu füllen sind (ISBN, DOI, Verlagsanschrift, Druckerei, Umschlag).
% Nur im Druck-PDF sichtbar; die Online-Version hat kein Impressum.
% ------------------------------------------------------------------
\lowertitleback{%
  \footnotesize
  \textbf{Roman Bartnik}\\
  Generative KI als Hilfe zum Lehren und Lernen\\[0.6em]
  Technische Hochschule Köln\\
  Fakultät für Informatik und Ingenieurwissenschaften, Campus Gummersbach\\[1em]
  Bibliografische Information der Deutschen Nationalbibliothek: Die Deutsche
  Nationalbibliothek verzeichnet diese Publikation in der Deutschen
  Nationalbibliografie; detaillierte bibliografische Daten sind im Internet
  über \url{https://dnb.de} abrufbar.\\[1em]
  Dieses Werk ist lizenziert unter einer Creative-Commons-Lizenz Namensnennung --
  Weitergabe unter gleichen Bedingungen 4.0 International (CC BY-SA 4.0):
  \url{https://creativecommons.org/licenses/by-sa/4.0/deed.de}. Ausgenommen sind
  Abbildungen und Zitate, die mit einem abweichenden Rechtehinweis gekennzeichnet sind.\\[1em]
  Gefördert aus dem Open-Access-Publikationsfonds der TH Köln.\\
  % \includegraphics[height=12mm]{images/logos/th-koeln-oa-fonds.png} % Logo liegt noch nicht vor
  {[}Logo des Publikationsfonds{]}\\[1em]
  Logos Verlag Berlin GmbH\\
  {[}Anschrift des Verlags{]}\\
  \url{https://www.logos-verlag.de}\\[1em]
  Erste Auflage, [Monat] 2026\\
  ISBN [978-3-xxxxx-xxx-x]\\
  DOI [10.30819/xxxx]\\[1em]
  Online-Version mit interaktiven Elementen und laufend aktualisierten Seiten:\\
  \url{https://th-koln-bartnik.github.io/genai4teaching/}\\[1em]
  Satz: Roman Bartnik mit Quarto und XeLaTeX\\
  Umschlag: [Gestaltung, verwendete Abbildung und Bildnachweis]\\
  Druck: [Druckerei]\\[1em]
  \textbf{Zitiervorschlag:} Bartnik, R. (2026). \emph{Generative KI als Hilfe zum
  Lehren und Lernen}. Logos Verlag Berlin. [DOI]%
}

% Kolumnentitel: lange Kapitel- und Abschnittstitel werden mit Auslassungspunkten gekürzt,
% statt über den Satzspiegel hinauszulaufen (KOMA-Script setzt links das Kapitel, rechts den Abschnitt).
\usepackage{truncate}
\renewcommand*{\TruncateMarker}{\,\textellipsis}
\renewcommand*{\chaptermark}[1]{\markboth{\MakeMarkcase{\truncate{0.92\textwidth}{\ifnumbered{chapter}{\chaptermarkformat}{}#1}}}{}}
\renewcommand*{\sectionmark}[1]{\markright{\MakeMarkcase{\truncate{0.92\textwidth}{\ifnumbered{section}{\sectionmarkformat}{}#1}}}}

